\documentclass{beamer}
\usepackage{graphics}

\title{APARAJIT: An Edge-compute System-on-chip}
\author{Madhav Desai\\
Department of Electrical Engineering\\
IIT Bombay}

\begin{document}
\maketitle

\frame[containsverbatim]{\frametitle{Overview}
\begin{itemize}
	\item  APARA, the edge compute accelerator.
		\begin{itemize}
			\item instruction set architecture.
			\item data and instruction memory.
			\item DMA engine.
		\end{itemize}
	\item  APARAJIT a  minimal system on chip:
		\begin{itemize}
			\item Control processor (AJIT processor).
			\item Accelerator (APARA accelerator).
			\item DMA controller.
			\item Switch.
			\item Memory.
		\end{itemize}
	\item Programming APARAJIT.
	\item Demonstration of FPGA prototype and early performance data.
	\item Road-map.
\end{itemize}
}

\frame[containsverbatim]{\frametitle{APARA architecture}

\begin{figure}
\includegraphics[width=8cm]{figs/APARAMT.pdf}
 \caption{APARA structure}
\end{figure}

}


\frame[containsverbatim]{\frametitle{APARA features}

\begin{itemize}
\item A memory-coupled accelerator.
\item Disjoint instruction and data memories.
\item Four execution threads multiplexed onto a single
	hardware pipeline.
\item Hardware pipeline: Bundles of instructions with up to 8 instructions per bundle.
\item Supported types: i4/i8/i16/i32/i64, f4/f8/f16/f32/f64
\item Per-thread 32x 64-bit registers, Per-thread program counters, Per-thread return address stack.
\item Instructions:
	\begin{itemize}
		\item Memory load/store:  from byte to 32-bytes.
		\item Scalar and vector instructions, including dot-product.
		\item Branching and subroutines.
	\end{itemize}
\end{itemize}
}

\frame[containsverbatim]{\frametitle{Multi-threaded APARA: functional description}
\begin{verbatim}
// 4 threads, 4x PC, 4x REGSETS, 4x RAS, 4x ERR
PC[1 to 4] = initvalues , REGS[1 to 4],
RAS[1 to 4] = empty, ERR[1 to 4] = 0, ACTIVE=1
while(1)
{ 
    bundle := fetch_bundle_from_imem(PC[ACTIVE])
    ops           := 
	   get_ops_needed_by_bundle(bundle, REGS[ACTIVE])
    out_values    := 
            execute_instructions_in_bundle(bundle, ops)
    PC[ACTIVE], REGS[ACTIVE], ERR[ACTIVE] := 
            update_registers_pc_ras_error (out_values)
    ACTIVE = next_active (ACTIVE, ERR[1 to 4]);    

    if(ACTIVE==0) break;
}
\end{verbatim}
}

\frame[containsverbatim]{\frametitle{Abstract machine: mcode program}
\begin{itemize}
	\item Up to 4 threads are active in the abstract machine.  They
		are multiplexed onto the execution pipeline.
	\item Each thread has its own state: program counter, 
		register file, return address stack, error status.
	\item Each thread executes an Mcode program.
	\item A Mcode program is a sequence of instruction bundles.
	\item Each bundle can have 1/2/4/8 instructions.
	\item Instructions in a bundle are executed in 
	parallel using a synchronous fetch-execute-write sequence which 
	guarantees no race conditions.
	\item A single bundle  can have at most one control transfer
	instruction.
	\item Two instructions in a bundle cannot write to
	the same register.
\end{itemize}
}



\frame[containsverbatim]{\frametitle{Bundle}

\begin{verbatim}
label: 
  ||
      <instr 1>
      <instr 2>
      <instr 3>
      ...
      <instr 8>
  ;
\end{verbatim}
Can have 1/2/4/8 instructions in each bundle.

}

\frame[containsverbatim]{\frametitle{Instruction set: instructions in a bundle}
\begin{itemize}
\item Each instruction is 32-bits wide.
\item Instruction classes:
	\begin{itemize}
		\item ALU instructions.
			\begin{itemize}
				\item Integer: $+,-,*,/,or,and,xor,nor,nand,xnor, <<, >>$.
				\item Float:   $+,-,*,/$, square-root.
			\end{itemize}
		\item Load/Store insrtuctions.
		\item Control transfer instructions: branch, call, return, conditional move.
		\item Vector instructions: vector add, subtract, multiply, scale, reduce, dot-product.
		\item Miscellaneous instructions: set, slice, pack, null, nop, cast.
	\end{itemize}
\item Control transfer instructions: branch, call, return.
\end{itemize}
}

\frame[containsverbatim]{\frametitle{Types}
\begin{verbatim}
Integer (signed, unsigned, vector)
  signed:   ($i4) ($i8) ($i16) ($i32) ($i64) 
  unsigned: ($u4) ($u8) ($u16) ($u32) ($u64) 
            ($u128) ($u256)
  float
            ($f4) ($f8) ($f16) ($f32) ($f64)

  vector:   ($vi4) ($vi8) ($vi16) ($vi32) 
            ($vu4) ($vu8) ($vu16) ($vu32) 
            ($vf4) ($vf8) ($vf16) ($vf32) 
\end{verbatim}

Note: 128 and 256 bit types are used only during data
	transfer with memory, not for arithmetic.

}

\frame[containsverbatim]{\frametitle{Orthogonalization:  op-code and type}
\begin{verbatim}
+ $r4 ($i16) $r1 $r2
+ $r4 ($i32) $r1 $r2 

+ $r4 ($i64) $r1 256

+ $r4 ($u16) $r1 $r2
+ $r4 ($u32) $r1 $r2 
\end{verbatim}
}

\frame[containsverbatim]{\frametitle{Vector ALU instructions}

\begin{verbatim}
$v * $r8  ($vi32) $r8  $r12
\end{verbatim}

Same as ALU instructions except that 
\begin{itemize}
	\item immediate operands not possible.
	\item source and destination types are vectors.
	\item Replicate flag is provided instead of immediates.
\end{itemize}
}

\frame[containsverbatim]{\frametitle{Vector ALU DOT instruction}
\begin{verbatim}
$dot                $r24 ($vi16) $r8  $r12
$dot $accumulate    $r24 ($vi16) $r8  $r12
\end{verbatim}

A dot-product on the two source vectors which
produces a destination value of the specified type ($rd = rs1.rs2$).  
The accumulate version does $rd = rs1.rs2 + rd$.

}

\frame[containsverbatim]{\frametitle{Vector VREDUCE instruction}

\begin{verbatim}
$vreduce + $r8  ($vi32) $r12
\end{verbatim}
The register \$rs1 is a vector.   The elements
of the source vector are added using precision specified
by dest-type, and the result is saved into \$rd.

}


\frame[containsverbatim]{\frametitle{Casts: the CAST instruction}

\begin{verbatim}
$cast ($f32)  $r8 ($f16) $r7
$cast ($vf32) $r8 ($vf16) $r7
$cast ($vf8)  $r8 ($vf16) $r9
\end{verbatim}

The bits in rs  are mapped to bits in rd via a cast
operation.   For scalars, it is the normal cast (as
in the first example).  For vectors, as in the
second example, the vector in \$r7  is mapped to
\$r8 and \$r9.  Two output registers are used.
In the third example, the destination vector is
filled with 8 values and this needs us to read
a multiplicity of two registers \$r9 and \$r10.

The multiplicity of in or out registers
should be at most 2.

}

\frame[containsverbatim]{\frametitle{Casts: the PACK instruction}
\begin{verbatim}
$pack $r8 64 32 $r4
$pack $r8 32 16 $r4
\end{verbatim}

In the first example, The 32 bit values in \$r4, \$r5 are concatenated
to produce a 64-bit result value in \$r8.   The concatenated values
must fit into a single destination register (in this case \$r8).
The second example is similar but two 16-bit values are packed
into a 32-bit value.

}

\frame[containsverbatim]{\frametitle{Data transfer: the LOAD instruction}
\begin{verbatim}
$ld ($i64)  $r3 [$r1 + $r0]
$ld ($u256) $r4 [$r1 + $r2 ++1]
$ld ($u64)  $r8 [$r1 + 256]
\end{verbatim}


The specified number of bits are fetched from memory
and kept in registers starting at \$rd.  Up to 64 bits
use one destination register. Higher values of number
of bits (max value is 256) use two to four destination
registers.   Auto-increment, if specified increments
rs2 by the specified value  times the size of transfer
in bytes.
}


\frame[containsverbatim]{\frametitle{Data transfer: the STORE instruction}
\begin{verbatim}
$st <dest-type> [ <rs1> + <rs2> <auto-incr>?] <rd>  OR
$st <dest-type> [ <rs1> + 16] <rd>  

$st ($i32)  [$r5 + $r6] $r8
$st ($i128) [$r5 + $r6] $r8
\end{verbatim}

The specified number (in dest-type) of bits are moved from
the specified rd register into memory.  Note that multiple
source registers may be used (e.g. for 128 bits we need
r8,r9 to be written to memory).  Register indices and
memory addresses must be aligned values.

}

\frame[containsverbatim]{\frametitle{Conditional move: CMOV}

\begin{verbatim}
$cmov ($uint32) $r5 > --1  $r6 $r7
\end{verbatim}
In the example, if \$r5 is greater than 0, then
\$r7 is moved to \$r6.   The \$r5 register value
is decremented.

}


\frame[containsverbatim]{\frametitle{Control Transfer: the BRANCH instruction}
\begin{verbatim}
? ($i32) $r5 >  --1  $goto loop_begin
\end{verbatim}

The example instruction above can be interpreted as follows:  Register \$r5 is compared
with $0$ and if it is $> 0$ (as a 32-bit signed integer) then control is transferred
to the label loop\_begin.  The value of \$r5 is decremented by $1$ after the comparison.
The auto-decrement (or increment) is optional.

}

\frame[containsverbatim]{\frametitle{Control Transfer: the CALL instruction}
\begin{verbatim}
$call foo 
$call $r5
\end{verbatim}
The expected next PC is saved on the return address stack and
control is passed to the bundle with label foo or
with the address $(PC + \$r5)$ (here \$r5 is interpreted
as a 32-bit signed integer).

}

\frame[containsverbatim]{\frametitle{Control Transfer: the RETURN instruction}
\begin{verbatim}
$return
\end{verbatim}

On a return instruction, the return address stack is popped
and control is passed to the popped address.
}

\frame[containsverbatim]{\frametitle{Miscellaneous: the SET instruction}
\begin{verbatim}
$set $r6  2 _hffff
\end{verbatim}

The 16-bit init-value $0xffff$ is inserted into r6
at bits [31:16].  SET is meant to used for initialization,
16 bits at a time.

}

\frame[containsverbatim]{\frametitle{Miscellaneous: the NULL, NOP and HALT instructions}
\begin{verbatim}
$null
$nop
$halt
\end{verbatim}

The NULL and NOP instructions do nothing.   The HALT instruction halts the
machine (after executing all instructions in the bundle in which the HALT occurs).

}

\frame[containsverbatim]{\frametitle{Tools}
\begin{itemize}
\item mcode\_align:  insert null instructions to align
bundles to required boundaries.
\item mcode\_assemble:  generate hex instructions for 
instruction memory.
\item mcode\_disassemble: from hex back to assembly.
\item mcode\_run:  an emulator.
\end{itemize}
}

\frame[containsverbatim]{\frametitle{Code example: simple dot product}
SEE THE CODE.
}

\frame[containsverbatim]{\frametitle{APARAJIT 1P1E: first prototype}

\begin{figure}
\includegraphics[width=8cm]{figs/1P1E.V2.pdf}
	\caption{APARAJIT 1P1E prototype (with Ethernet and PCIe interfaces shown)}
\end{figure}

}

\frame[containsverbatim]{\frametitle{APARAJIT 1P1E Prototype: Notes}
\begin{itemize}
	\item Using the configuration interface, we can start up to four
		execution threads by supplying an initial PC to each of them.
	\item Using the configuration interface, we can start up to four
		DMA threads by supplying source and destination 
		parameters to the DMA thread.
	\item   LUT count for processor (with FPU, MMU, L1 caches): 80K LUT.
	\item   LUT count for accelerator engine:  35K LUT per lane.
		\begin{itemize}
			\item Each lane:  2x FP32 mul, add, FP64 mul, add, 64-bit integer add, mul,
				4/8/16 bit integer and FP dot product unit.
		\end{itemize}
\end{itemize}
}


\frame[containsverbatim]{\frametitle{APARAJIT 1P1E Prototype: Simple Application Flowchart}

\begin{figure}
\includegraphics[width=8cm]{figs/Flowchart.pdf}
 \caption{Flow-chart}
\end{figure}

}

\frame[containsverbatim]{\frametitle{APARAJIT 1P1E Prototype: peer-to-peer queueing model}

\begin{figure}
\includegraphics[width=8cm]{figs/Flowchart.pdf}
 \caption{Flow-chart}
\end{figure}

}

\frame[containsverbatim]{\frametitle{APARAJIT 1P1E Prototype: DMA slicing}

\begin{figure}
\includegraphics[width=8cm]{figs/DmaSlicing.pdf}
 \caption{DMA Slicing}
\end{figure}
}

\frame[containsverbatim]{\frametitle{APARAJIT 1P1E Prototype: DMA slicing demo}

\begin{itemize}
	\item Configure a main memory to data memory DMA transfer on a DMA thread.
	\item Start the DMA thread.
	\item Wait until the DMA thread is completed.
	\item Configure a data memory to main memory DMA thread.
	\item Start the DMA thread.
	\item Wait until the DMA thread is completed.
	\item Repeat for each thread.
\end{itemize}

}


\frame[containsverbatim]{\frametitle{APARAJIT 1P1E Prototype: DMA performance demo}

\begin{itemize}

	\item Configure a main memory to data memory DMA transfer on a DMA thread.
	\item Start the DMA thread.
	\item Wait until the DMA thread is completed.
	\item Configure a data memory to main memory DMA thread.
	\item Start the DMA thread.
	\item Wait until the DMA thread is completed.
	\item Repeat for each thread.
	\item Observed performance: 1.2K cycles to transfer 8KB from main memory to local memory.
	\item Observed performance: 1.1K cycles to transfer 8KB from local memory to main memory.

\end{itemize}

}

\frame[containsverbatim]{\frametitle{APARAJIT 1P1E Prototype: 32-bit integer DOT-product 2K x 2K  demo} 
\begin{itemize}

	\item Set up initial conditions of registers (source pointers of two arrays, number of elements).
	\item Start four exec threads: all do the same work.
	\item Wait until all four threads have completed.
	\item Dump registers.
	\item Measure timings (compare with processor: 31X speedup).

\end{itemize}

}

\frame[containsverbatim]{\frametitle{APARAJIT 1P1E Prototype: 8-bit DOT-product 8K x 8K  demo}
\begin{itemize}

	\item Set up initial conditions of registers (source pointers of two arrays, number of elements).
	\item Start four exec threads: all do the same work.
	\item Wait until all four threads have completed.
	\item Dump registers.
	\item Measure timings (compare with processor: 155X speedup).

\end{itemize}

}

\frame[containsverbatim]{\frametitle{APARAJIT 1P1E Prototype: queueing network based application architecture}

\begin{figure}
\includegraphics[width=8cm]{figs/Queueing.pdf}
 \caption{Queueing}
\end{figure}

}

\frame[containsverbatim]{\frametitle{Beyond APARAJIT 1P1E Prototype:2P2E, 4P2E, 4P4E etc.}
\begin{itemize}
	\item Add Ethernet, PCIE interfaces to SOC.
	\item mPnE: m processors, n engines.  
		\begin{itemize}
			\item 2P1E, 2P2E, 4P1E, 4P2E etc...
		\end{itemize}
	\item Variants: with 1/2/4 lanes per accelerator for deep embedded applications.
	\item Clusters to scale beyond 4P4E.
\end{itemize}
}

\frame[containsverbatim]{\frametitle{Software for APARAJIT SOC}
\begin{itemize}
	\item Bare metal drivers.
	\item RTOS port to SOC: micro-kernel based OS identified (Redkill, NuttX).
	\item Software library:  BLAS, FFT, AIML operations.
		\begin{itemize}
			\item Initially in assembly.
			\item Eventually: compile from high level language.
		\end{itemize}
\end{itemize}
}

\frame[containsverbatim]{\frametitle{Collaborations: APARAJIT SOC}
\begin{itemize}
	\item Accelerator kernel routines.
	\item Accelerator evaluation/use in systems.
		\begin{itemize}
			\item as building block.
		\end{itemize}
	\item Domain specific use cases: bio-medical, robotics.
	\item Other ideas?
\end{itemize}
}

\end{document}
